\chapter{实验总结}
\section{完成情况与证据链}
本实验完成Week 3--8的连续工程：底层能够创建、查询和流式追加不等长归并段；Project 1完成分块内排与后序二路计划；Project 2完成败者树置换选择、冻结解冻、三线程双缓冲与哈夫曼计划；Project 3完成共享$2k$输入池、动态预取、真正K路选择与多轮主控。

对机制的确认不只依赖最终排序结果。小例子给出段长、树节点、冻结标记和块号变化；独立断言检查一条路占用四块与跨路回收；220次端到端测试逐项对照完整参考排序；Benchmark每行也经过精确验证。Windows的九个具名配置、F5断点、F10单步、Ctrl+F5以及中文路径构建经过实际运行。源码注释最后独立提交，随后再次编译并通过Debug/Release完整验收。

\section{实际遇到的问题与解决过程}
\subsection{构建环境与编码差异}
起初Windows普通终端没有可调用的C++编译器，VS Code扩展本身也不提供编译器。先使用已存在的WSL GCC完成算法验证，后经用户确认安装微软Build Tools的必要C++组件，复用现有CMake。Windows使用Visual Studio多配置生成器，Debug和Release程序位于各自子目录；不能照搬Linux单配置目录中的程序路径。

PowerShell默认读取UTF-8中文材料曾显示乱码，改为显式UTF-8读取后恢复。原生MSVC编译保留\code{/utf-8}，中文脚本采用兼容Windows PowerShell的UTF-8 BOM，文档和源码保持UTF-8。解决办法是核对读写编码及实际编译日志，而非删除中文注释。此次报告中的源码快照仍保留所有中文说明。

\subsection{测试目录生命周期与文件系统开销}
此前WSL把大量微小I/O测试放在Windows挂载路径，Project 1阶段测试约97.7秒；改到Linux用户目录后，同类运行明显加快。因此原始WSL性能表明确使用Linux文件系统，不与挂载路径数据混合。另一次WSL空闲退出后\code{/tmp}构建目录消失，最终改用用户缓存目录，避免临时目录生命周期影响复现。

测试程序使用固定文件名，便于课堂演示和检查目录，但同一目录的多个进程会互相覆盖。当前一键验收顺序执行，Debug、Release及CTest使用明确的工作位置。目录隔离是实验条件的一部分，不应在无意并发运行后把文件冲突解释为归并算法错误。

\subsection{Ctrl+F5复用终端的工作目录}
实际GUI验收发现，C/C++扩展1.35.6在Ctrl+F5复用已有终端时未更新工作目录，可能将相对路径数据写到仓库根。F5调试子进程的目录正确，但终端Shell的提示符仍在原位置，单看提示符不能确认子进程的文件位置。

检查扩展实际行为后，采用最小的Windows入口兼容修正：具名启动配置传\code{--workdir DIRECTORY MODE}，程序先检查目录并切换，再沿用原模式解析。该修正只放在Windows条件编译分支，不改排序接口。中文目录成功、无效目录拒绝、缺参数拒绝均经测试；真实Ctrl+F5后运行目录文件更新时间改变，根目录旧文件保持不变。该修正与后续纯注释提交分开保存，避免把行为变化伪装成注释维护。

\subsection{接口连接、线程收尾和验证口径}
先前一次在旧构建仍运行时接入下一阶段源码，CMake源列表与引用短暂不同步，出现未链接\code{replacementSelection}；按完整列表重新配置构建后通过。最终回归期间先停止源码修改，再构建完整工程，避免用更新中的文件混合证明最终状态。

审查中为线程创建失败补充先join已启动线程的收尾，防止仍可连接的线程析构触发终止。没有故意制造线程资源耗尽，因此只陈述审查与正常运行验证。尾块、段边界和未知头值分别通过明确状态处理与独立测试覆盖；本报告不虚构曾发生的死锁或丢数事故。

验证口径也作了明确区分。自动生成输入能够完整对照\code{std::sort}；用户自定义大文件仅流式检查数量和非降序，并在输出中说明不是精确多重集合比较。保存的Benchmark计时不包含参考检查，不将“测试总耗时”当成排序时间，也不把应用层调用称为硬盘物理I/O。

\section{学习收获与验收理解}
本项目把排序中的抽象“序列”落实为文件中的连续区域。理解一个段不仅需要知道值的顺序，还要知道目录下标、段内元素位置和字节偏移的转换。\code{beginRun/writeBlock/endRun}使段长未知时仍能持续输出，缓冲区则让算法保持逐元素逻辑而文件保持分块访问。

归并计划与败者树解决的是两类不同问题。前者安排哪些段先合并，叶子是段；后者决定下一个输出哪个候选，叶子是当前值。哈夫曼计划优化加权搬移量，冠军路径重赛减少重复比较，置换选择通过冻结延长初始段。把它们混成一棵“排序树”会妨碍理解，按接口和生命周期区分后更容易解释复用关系。

线程部分最关键的是数据何时可见、由谁修改。预留块尚未入队，计算不能读取；耗尽块尚未同步，输入不能重用；新输出尚未完成，写线程不能访问。join不是装饰性等待，而是这些状态交接的边界。实际时间又说明，逻辑上存在并行任务，不意味着在小数据、小块环境下一定更快。

学习和验收时可沿同一15值输入追踪三个项目：先定位R0的物理偏移，再画普通计划，随后手算冻结与解冻，最后指出物理块2怎样跨路使用。报告与中文注释提供可检查依据；对函数的理解仍应通过自己预测下一步状态、修改参数并核对实际结果来巩固。

\section{局限性与改进方向}
\begin{enumerate}
\item 文件保留全部初始和中间段，便于观察但增加磁盘占用；尚无空间回收和压缩。
\item 文件结构采用本机字节序，未实现通用序列化、校验和与掉电恢复。输入输出须为不同真实文件，未处理所有路径别名或链接。
\item 候选和数据块受容量约束，目录、段号列表及计划仍为$O(R)$；并非恒定总内存系统。
\item 每阶段新建三个线程，未实现线程池或复杂异步调度。可考虑复用线程、扩大块并重复测量，但这些不属于已完成结果。
\item 未开展大于物理内存的长期磁盘压力测试，未验证磁盘满、真实掉电、跨字节序及系统线程耗尽。测试不能证明所有未覆盖故障路径。
\item K路每阶段扫描与重建树、段名线性查询及重复加载目录存在额外成本；后续若优化，应先保持机制与正确性测试，再单独测量收益。
\end{enumerate}

\section{材料来源与复现文件}
课程依据为用户提供的《外存数据管理》85页课件及两份课堂语音转写；六部分报告结构采用用户提供的Project Reports要求。原课件与讲稿仅供本地参考，未发布到公开仓库。本报告的架构图、算法图和说明按最终源码重新绘制和撰写。

实现依据为最终仓库17个C++文件、\code{README.md}及\code{docs/}中的讲义、课程对应、实验总结与验收问答。Windows与WSL数据分别来自\code{benchmark-windows.csv}、\code{benchmark.csv}；截图为真实Windows程序窗口。最后回归原始日志保存在本地\code{.local/windows/}，报告目录另外保留必要证据副本。

\code{report/README.md}给出XeLaTeX编译方法、源码快照生成方法和检查步骤。\code{scripts/build_report.ps1}执行编译与验证，\code{scripts/make_charts.py}从CSV生成矢量图。源代码附录紧接本章并位于整份报告末尾，包含最终实际实现及中文注释。
